%%%%%%%%%%%%%%%%
%%% deut 4 %%%
%%%%%%%%%%%%%%%%

\Chapter{4}

\Verse{1}
{
	\hDtr1{וְעַתָּה יִשְׂרָאֵל שְׁמַע אֶל הַחֻקִּים וְאֶל הַמִּשְׁפָּטִים אֲשֶׁר אָנֹכִי מְלַמֵּד אֶתְכֶם לַעֲשׂוֹת לְמַעַן תִּחְיוּ וּבָאתֶם וִירִשְׁתֶּם אֶת הָאָרֶץ אֲשֶׁר יְהֹוָה אֱלֹהֵי אֲבֹתֵיכֶם נֹתֵן לָכֶם׃}
}
{
	\eDtr1{
		4 “And now, Israel, listen to the laws and to the
		judgments that I’m teaching you to do, so that you’ll
		live, and you’ll come and take possession of the land that
		YHWH, your fathers’ God, is giving you
	}
}
{
%	\footnote{
%		so that you’ll live. This theme of choosing life begins here and will
%		recur several times in the course of Moses’ speech, and culminate in his
%		last sentences (Deut 30:19). See the comment there.
%	}
}

\Verse{2}
{
	\hDtr1{לֹא תֹסִפוּ עַל הַדָּבָר אֲשֶׁר אָנֹכִי מְצַוֶּה אֶתְכֶם וְלֹא תִגְרְעוּ מִמֶּנּוּ לִשְׁמֹר אֶת מִצְוֺת יְהֹוָה אֱלֹהֵיכֶם אֲשֶׁר אָנֹכִי מְצַוֶּה אֶתְכֶם׃}
}
{
	\eDtr1{
		You shall not add onto the thing that I command you, and
		you shall not subtract from it: observing the commandments
		of YHWH, your God, that I command you.
	}
}
{
%	\footnote{
%		You shall not add … and you shall not subtract. This is an essential
%		command: it governs the way Israel is to observe all the other
%		commandments. The second part—not to subtract from the law—is more
%		obvious than the first, but that makes the recognition of the first part
%		all the more important. One may think that, by doing more than the law
%		requires, one is doing better, being more religious, more observant,
%		when one is in fact thus violating the law. And when one imposes such
%		additions to the law on others, one puts them at risk of violation as
%		well. This is stated here in the last book of the Torah as a command,
%		but it was already conveyed through a story at the beginning of the
%		Torah, when the woman says to the snake that God commanded that they not
%		eat from the tree or touch it, though God had said nothing about
%		touching it (see the comment on Gen 3:3). Adding to a command is as
%		dangerous as taking away from it. Indeed, as Rashi noted in his comment
%		on that verse in Genesis, adding to it may lead to taking away from it.
%		One must, therefore, use extreme caution in interpreting and expanding
%		the laws in the Torah. This is sometimes a delicate task because in
%		postbiblical Judaism a principle developed of “building a fence around
%		the Torah.” One starts to observe the Sabbath early, before sunset, so
%		as not to risk violating the Sabbath by mistakenly starting late. But
%		there is a difference between sincere care for the law, which is done
%		out of reverence for it, and extreme expansion of the law, which is done
%		for other motives: fear, power over others. The law is presented in the
%		Torah as divine, but it requires human strength and wisdom to carry it
%		out.
%	}
}

\Verse{3}
{
	\hDtr1{עֵינֵיכֶם הָרֹאוֹת אֵת אֲשֶׁר עָשָׂה יְהֹוָה בְּבַעַל פְּעוֹר כִּי כל הָאִישׁ אֲשֶׁר הָלַךְ אַחֲרֵי בַעַל פְּעוֹר הִשְׁמִידוֹ יְהֹוָה אֱלֹהֶיךָ מִקִּרְבֶּךָ׃}
}
{
	\eDtr1{
		It’s your eyes that saw what YHWH did at Baal Peor, that
		every man who went after Baal Peor: YHWH, your God,
		destroyed him from among you.
	}
}
{
}

\Verse{4}
{
	\hDtr1{וְאַתֶּם הַדְּבֵקִים בַּיהֹוָה אֱלֹהֵיכֶם חַיִּים כֻּלְּכֶם הַיּוֹם׃}
}
{
	\eDtr1{
		But you who were clinging to YHWH, your God: you’re all
		alive today!
	}
}
{
}

\Verse{5}
{
	\hDtr1{רְאֵה לִמַּדְתִּי אֶתְכֶם חֻקִּים וּמִשְׁפָּטִים כַּאֲשֶׁר צִוַּנִי יְהֹוָה אֱלֹהָי לַעֲשׂוֹת כֵּן בְּקֶרֶב הָאָרֶץ אֲשֶׁר אַתֶּם בָּאִים שָׁמָּה לְרִשְׁתָּהּ׃}
}
{
	\eDtr1{
		See: I’ve taught you laws and judgments as YHWH, my God,
		commanded me, to do so within the land to which you’re
		coming to take possession of it.
	}
}
{
}

\Verse{6}
{
	\hDtr1{וּשְׁמַרְתֶּם וַעֲשִׂיתֶם כִּי הִוא חכְמַתְכֶם וּבִינַתְכֶם לְעֵינֵי הָעַמִּים אֲשֶׁר יִשְׁמְעוּן אֵת כּל הַחֻקִּים הָאֵלֶּה וְאָמְרוּ רַק עַם חָכָם וְנָבוֹן הַגּוֹי הַגָּדוֹל הַזֶּה׃}
}
{
	\eDtr1{
		And you shall observe and do them because it’s your wisdom
		and your understanding before the eyes of the peoples, in
		that they will hear all these laws and will say, ‘Only a
		wise and understanding people is this great nation.’
	}
}
{
}

\Verse{7}
{
	\hDtr1{כִּי מִי גוֹי גָּדוֹל אֲשֶׁר לוֹ אֱלֹהִים קְרֹבִים אֵלָיו כַּיהֹוָה אֱלֹהֵינוּ בְּכל קרְאֵנוּ אֵלָיו׃}
}
{
	\eDtr1{
		Because who is a great nation that has gods close to it
		like YHWH, our God, in all our calling to him?
	}
}
{
}

\Verse{8}
{
	\hDtr1{וּמִי גּוֹי גָּדוֹל אֲשֶׁר לוֹ חֻקִּים וּמִשְׁפָּטִים צַדִּיקִם כְּכֹל הַתּוֹרָה הַזֹּאת אֲשֶׁר אָנֹכִי נֹתֵן לִפְנֵיכֶם הַיּוֹם׃}
}
{
	\eDtr1{
		And who is a great nation that has just laws and judgments
		like all this instruction that I’m putting in front of you
		today?
	}
}
{
}

\Verse{9}
{
	\hDtr1{רַק הִשָּׁמֶר לְךָ וּשְׁמֹר נַפְשְׁךָ מְאֹד פֶּן תִּשְׁכַּח אֶת הַדְּבָרִים אֲשֶׁר רָאוּ עֵינֶיךָ וּפֶן יָסוּרוּ מִלְּבָבְךָ כֹּל יְמֵי חַיֶּיךָ וְהוֹדַעְתָּם לְבָנֶיךָ וְלִבְנֵי בָנֶיךָ׃}
}
{
	\eDtr1{
		“Only be watchful, and watch yourself very much, in case
		you’ll forget the things that your eyes have seen and in
		case they’ll turn away from your heart, all the days of
		your life. But you shall make them known to your children
		and to your children’s children:
	}
}
{
}

\Verse{10}
{
	\hDtr1{יוֹם אֲשֶׁר עָמַדְתָּ לִפְנֵי יְהֹוָה אֱלֹהֶיךָ בְּחֹרֵב בֶּאֱמֹר יְהֹוָה אֵלַי הַקְהֶל לִי אֶת הָעָם וְאַשְׁמִעֵם אֶת דְּבָרָי אֲשֶׁר יִלְמְדוּן לְיִרְאָה אֹתִי כּל הַיָּמִים אֲשֶׁר הֵם חַיִּים עַל הָאֲדָמָה וְאֶת בְּנֵיהֶם יְלַמֵּדוּן׃}
}
{
	\eDtr1{
		the day that you stood in front of YHWH, your God, at
		Horeb, when YHWH said to me, ‘Assemble the people to me,
		and I’ll make them hear my words so that they’ll learn to
		fear me all the days that they’re living on the land, and
		they’ll teach their children.’
	}
}
{
}

\Verse{11}
{
	\hDtr1{וַתִּקְרְבוּן וַתַּעַמְדוּן תַּחַת הָהָר וְהָהָר בֹּעֵר בָּאֵשׁ עַד לֵב הַשָּׁמַיִם חֹשֶׁךְ עָנָן וַעֲרָפֶל׃}
}
{
	\eDtr1{
		And you came forward and stood below the mountain, and the
		mountain was burning in fire to the heart of the skies:
		darkness, cloud, and nimbus.
	}
}
{
%	\footnote{
%		the mountain was burning. The text in Exodus does not in fact use this
%		word for burning. It says that “Mount Sinai was all smoke because YHWH
%		came down on it in fire” (Exod 19:18). This word was rather the term for
%		the burning bush. And so we learn only now that the burning bush that
%		Moses experienced was already a sign of what was to come. And this in
%		turn may answer the old question of what is meant when God says to Moses
%		at the bush, “This is the sign for you that I have sent you. When you
%		bring out the people from Egypt you shall serve God on this mountain”
%		(Exod 3:12). What is the sign? The miraculous burning bush itself. And
%		it prefigures what will happen when the people come back to that
%		mountain. This point is supported later in this chapter (4:36) when
%		Moses reminds the people that at Sinai “you heard His voice from inside
%		the fire”—which recalls how Moses himself heard God call “from inside
%		the bush” that was on fire (Exod 3:4).
%	}
}

\Verse{12}
{
	\hDtr1{וַיְדַבֵּר יְהֹוָה אֲלֵיכֶם מִתּוֹךְ הָאֵשׁ קוֹל דְּבָרִים אַתֶּם שֹׁמְעִים וּתְמוּנָה אֵינְכֶם רֹאִים זוּלָתִי קוֹל׃}
}
{
	\eDtr1{
		And YHWH spoke to you from inside the fire. You were
		hearing the sound of words, but you weren’t seeing a form,
		just sound.
	}
}
{
}

\Verse{13}
{
	\hDtr1{וַיַּגֵּד לָכֶם אֶת בְּרִיתוֹ אֲשֶׁר צִוָּה אֶתְכֶם לַעֲשׂוֹת עֲשֶׂרֶת הַדְּבָרִים וַיִּכְתְּבֵם עַל שְׁנֵי לֻחוֹת אֲבָנִים׃}
}
{
	\eDtr1{
		And He told you His covenant that He commanded you to do,
		the Ten commandments, and He wrote them on two tablets of
		stones.
	}
}
{
%	\footnote{
%		His covenant … the Ten Commandments. The Ten commandments are the
%		covenant. All the other commandments are important, and there are
%		rewards for keeping them and consequences for breaking them. But the Ten
%		Commandments on the tablets are different. Violation of them by the
%		community (for example, widespread pagan worship or injustice) risks the
%		breaking of the covenant. They are the essence of the covenant itself.
%	}
}

\Verse{14}
{
	\hDtr1{וְאֹתִי צִוָּה יְהֹוָה בָּעֵת הַהִוא לְלַמֵּד אֶתְכֶם חֻקִּים וּמִשְׁפָּטִים לַעֲשֹׂתְכֶם אֹתָם בָּאָרֶץ אֲשֶׁר אַתֶּם עֹבְרִים שָׁמָּה לְרִשְׁתָּהּ׃}
}
{
	\eDtr1{
		And YHWH commanded me at that time to teach you laws and
		judgments, for you to do them in the land to which you’re
		crossing to take possession of it.
	}
}
{
}

\Verse{15}
{
	\hDtr1{וְנִשְׁמַרְתֶּם מְאֹד לְנַפְשֹׁתֵיכֶם כִּי לֹא רְאִיתֶם כּל תְּמוּנָה בְּיוֹם דִּבֶּר יְהֹוָה אֲלֵיכֶם בְּחֹרֵב מִתּוֹךְ הָאֵשׁ׃}
}
{
	\eDtr1{
		And you shall be very watchful of yourselves—because you
		didn’t see any form in the day that YHWH spoke to you at
		Horeb from inside the fire—
	}
}
{
}

\Verse{16}
{
	\hDtr1{פֶּן תַּשְׁחִתוּן וַעֲשִׂיתֶם לָכֶם פֶּסֶל תְּמוּנַת כּל סָמֶל תַּבְנִית זָכָר אוֹ נְקֵבָה׃}
}
{
	\eDtr1{
		in case you’ll be corrupted, and you’ll make a statue, a
		form of any figure, a design of a male or a female,
	}
}
{
}

\Verse{17}
{
	\hDtr1{תַּבְנִית כּל בְּהֵמָה אֲשֶׁר בָּאָרֶץ תַּבְנִית כּל צִפּוֹר כָּנָף אֲשֶׁר תָּעוּף בַּשָּׁמָיִם׃}
}
{
	\eDtr1{
		a design of any animal that’s in the earth, a design of
		any winged bird that flies in the skies,
	}
}
{
}

\Verse{18}
{
	\hDtr1{תַּבְנִית כּל רֹמֵשׂ בָּאֲדָמָה תַּבְנִית כּל דָּגָה אֲשֶׁר בַּמַּיִם מִתַּחַת לָאָרֶץ׃}
}
{
	\eDtr1{
		a design of anything that creeps on the ground, a design
		of any fish that’s in the water under the earth,
	}
}
{
}

\Verse{19}
{
	\hDtr1{וּפֶן תִּשָּׂא עֵינֶיךָ הַשָּׁמַיְמָה וְרָאִיתָ אֶת הַשֶּׁמֶשׁ וְאֶת הַיָּרֵחַ וְאֶת הַכּוֹכָבִים כֹּל צְבָא הַשָּׁמַיִם וְנִדַּחְתָּ וְהִשְׁתַּחֲוִיתָ לָהֶם וַעֲבַדְתָּם אֲשֶׁר חָלַק יְהֹוָה אֱלֹהֶיךָ אֹתָם לְכֹל הָעַמִּים תַּחַת כּל הַשָּׁמָיִם׃}
}
{
	\eDtr1{
		and in case you’ll raise your eyes to the skies, and
		you’ll see the sun and the moon and the stars, all the
		array of the skies, and you’ll be moved so that you’ll bow
		to them and serve them, when YHWH, your God, has allocated
		them to all the peoples under all the skies.
	}
}
{
%	\footnote{
%		the array of the skies. Again the last book of the Torah recalls the
%		beginning of the Torah. The term “array” has not been used to refer to
%		the heavenly bodies since the creation story (Gen 2:1). There it was
%		used in the account of God’s creation of the cosmos. Now Moses warns
%		that one might be inspired by their awesome quality to worship them.
%		Footnote 4:19. allocated them to all the peoples. This is doubly
%		extraordinary: (1) God is understood to have provided the pagan deities
%		that the other nations worship. (2) There is no criticism here of the
%		other nations for worshiping them. There is nothing wrong with other
%		people’s following their own faiths. They are only criticized elsewhere
%		in the Torah for specific practices in their religions that are
%		abhorrent to Israelites, such as human sacrifice.
%	}
}

\Verse{20}
{
	\hDtr1{וְאֶתְכֶם לָקַח יְהֹוָה וַיּוֹצִא אֶתְכֶם מִכּוּר הַבַּרְזֶל מִמִּצְרָיִם לִהְיוֹת לוֹ לְעַם נַחֲלָה כַּיּוֹם הַזֶּה׃}
}
{
	\eDtr1{
		But YHWH has taken you, and He has brought you out from
		the iron furnace, from Egypt, to become a legacy people
		for him, as it is this day.
	}
}
{
}

\Verse{21}
{
	\hDtr1{וַיהֹוָה הִתְאַנַּף בִּי עַל דִּבְרֵיכֶם וַיִּשָּׁבַע לְבִלְתִּי עבְרִי אֶת הַיַּרְדֵּן וּלְבִלְתִּי בֹא אֶל הָאָרֶץ הַטּוֹבָה אֲשֶׁר יְהֹוָה אֱלֹהֶיךָ נֹתֵן לְךָ נַחֲלָה׃}
}
{
	\eDtr1{
		“And YHWH had been incensed at me over your matters, and
		He swore that I would not cross the Jordan and not come to
		the good land that YHWH, your God, is giving you as a
		legacy.
	}
}
{
}

\Verse{22}
{
	\hDtr1{כִּי אָנֹכִי מֵת בָּאָרֶץ הַזֹּאת אֵינֶנִּי עֹבֵר אֶת הַיַּרְדֵּן וְאַתֶּם עֹבְרִים וִירִשְׁתֶּם אֶת הָאָרֶץ הַטּוֹבָה הַזֹּאת׃}
}
{
	\eDtr1{
		So I am dying in this land. I’m not crossing the Jordan.
		But you are crossing, and you shall possess this good
		land.
	}
}
{
%	\footnote{
%		good. The word “good” recurs nine times in the first chapters of
%		Deuteronomy. It culminates when Moses ends this section of his speech
%		with the instruction to keep all the commandments “so it will be good
%		for you” (4:40). It is reminiscent of Genesis 1. There God declares each
%		stage of the creation to be “good.” That is, the world starts off in a
%		state that is positive, but it soon goes into a variety of troubles
%		relating to human behavior (“The ground is cursed on your account,” Gen
%		3:17). And now Moses tells the people over and over that the land they
%		are entering is “good.” It too starts off in a positive state, and they
%		have the opportunity to keep it good or to ruin it.
%	}
}

\Verse{23}
{
	\hDtr1{הִשָּׁמְרוּ לָכֶם פֶּן תִּשְׁכְּחוּ אֶת בְּרִית יְהֹוָה אֱלֹהֵיכֶם אֲשֶׁר כָּרַת עִמָּכֶם וַעֲשִׂיתֶם לָכֶם פֶּסֶל תְּמוּנַת כֹּל אֲשֶׁר צִוְּךָ יְהֹוָה אֱלֹהֶיךָ׃}
}
{
	\eDtr1{
		Watch yourselves, in case you’ll forget the covenant of
		YHWH, your God, that He made with you, and you’ll make a
		statue of any form about which YHWH, your God, has
		commanded you,
	}
}
{
}

\Verse{24}
{
	\hDtr1{כִּי יְהֹוָה אֱלֹהֶיךָ אֵשׁ אֹכְלָה הוּא אֵל קַנָּא׃}
}
{
	\eDtr1{
		because YHWH, your God: He is a consuming fire, a jealous
		God.
	}
}
{
%	\footnote{
%		a consuming fire, a jealous God. Here Moses emphasizes the harsh,
%		frightening side of God. But he immediately softens it in his next few
%		sentences. See the comments on vv. 29 and 31 below. (On the term
%		“jealous,” see the comment on Exod 20:5.)
%	}
}

\Verse{25}
{
	\hDtr2{כִּי תוֹלִיד בָּנִים וּבְנֵי בָנִים וְנוֹשַׁנְתֶּם בָּאָרֶץ וְהִשְׁחַתֶּם וַעֲשִׂיתֶם פֶּסֶל תְּמוּנַת כֹּל וַעֲשִׂיתֶם הָרַע בְּעֵינֵי יְהֹוָה אֱלֹהֶיךָ לְהַכְעִיסוֹ׃}
}
{
	\eDtr2{
		“When you’ll produce children and grand-children and will
		have been in the land long, and you’ll be corrupt and make
		a statue of any form and do what is bad in the eyes of
		YHWH, your God, to provoke Him,
	}
}
{
}

\Verse{26}
{
	\hDtr2{הַעִידֹתִי בָכֶם הַיּוֹם אֶת הַשָּׁמַיִם וְאֶת הָאָרֶץ כִּי אָבֹד תֹּאבֵדוּן מַהֵר מֵעַל הָאָרֶץ אֲשֶׁר אַתֶּם עֹבְרִים אֶת הַיַּרְדֵּן שָׁמָּה לְרִשְׁתָּהּ לֹא תַאֲרִיכֻן יָמִים עָלֶיהָ כִּי הִשָּׁמֵד תִּשָּׁמֵדוּן׃}
}
{
	\eDtr2{
		I call the skies and the earth to witness regarding you
		today that you’ll perish quickly from the land to which
		you’re crossing the Jordan to take possession of it. You
		won’t extend days on it, but you’ll be destroyed!
	}
}
{
}

\Verse{27}
{
	\hDtr2{וְהֵפִיץ יְהֹוָה אֶתְכֶם בָּעַמִּים וְנִשְׁאַרְתֶּם מְתֵי מִסְפָּר בַּגּוֹיִם אֲשֶׁר יְנַהֵג יְהֹוָה אֶתְכֶם שָׁמָּה׃}
}
{
	\eDtr2{
		And YHWH will scatter you among the peoples, and you’ll be
		left few in number among the nations where YHWH will drive
		you.
	}
}
{
%	\footnote{
%		scatter. The threat of scattering a people is frightening itself, and,
%		again, it harks back to the early part of Genesis, where YHWH scatters
%		the people of the world in the matter of the tower of Babylon (Gen
%		11:9). Footnote 4:27. few in number. The other occurrence of this
%		expression in the Torah is where Jacob reprimands his sons Simeon and
%		Levi for massacring the city of Shechem; he says that he is vulnerable
%		now because he is few in number (Gen 34:30). And the denouement of this
%		matter comes when, on his deathbed, Jacob condemns Simeon and Levi to be
%		scattered (Gen 49:7), using the same word that is used for scattering in
%		the first half of this verse. The double parallel makes a notable
%		reminiscence of Genesis here.
%	}
}

\Verse{28}
{
	\hDtr2{וַעֲבַדְתֶּם שָׁם אֱלֹהִים מַעֲשֵׂה יְדֵי אָדָם עֵץ וָאֶבֶן אֲשֶׁר לֹא יִרְאוּן וְלֹא יִשְׁמְעוּן וְלֹא יֹאכְלוּן וְלֹא יְרִיחֻן׃}
}
{
	\eDtr2{
		And you’ll serve gods, the work of human hands, wood and
		stone, there, that don’t see and don’t hear and don’t eat
		and don’t smell.
	}
}
{
%	\footnote{
%		work of human hands. This verse is a brief expression of what is said in
%		Ps 115:4–8.
%	}
}

\Verse{29}
{
	\hDtr2{וּבִקַּשְׁתֶּם מִשָּׁם אֶת יְהֹוָה אֱלֹהֶיךָ וּמָצָאתָ כִּי תִדְרְשֶׁנּוּ בְּכל לְבָבְךָ וּבְכל נַפְשֶׁךָ׃}
}
{
	\eDtr2{
		“But if you’ll seek YHWH, your God, from there, then
		you’ll find Him, when you’ll inquire of Him with all your
		heart and all your soul.
	}
}
{
%	\footnote{
%		But if you’ll seek … then you’ll find. Even though Moses has just told
%		the people that God will destroy and scatter them, he still returns to
%		the idea that they can repent and that God will forgive them. The
%		conception of God is always based in compassion. No matter what the
%		people’s offense, God is forgiving and merciful, and there can be a new
%		reconciliation.
%	}
}

\Verse{30}
{
	\hDtr2{בַּצַּר לְךָ וּמְצָאוּךָ כֹּל הַדְּבָרִים הָאֵלֶּה בְּאַחֲרִית הַיָּמִים וְשַׁבְתָּ עַד יְהֹוָה אֱלֹהֶיךָ וְשָׁמַעְתָּ בְּקֹלוֹ׃}
}
{
	\eDtr2{
		When you have trouble, and all these things have found
		you, in the future days, if you’ll go back to YHWH, your
		God, and listen to His voice,
	}
}
{
}

\Verse{31}
{
	\hDtr2{כִּי אֵל רַחוּם יְהֹוָה אֱלֹהֶיךָ לֹא יַרְפְּךָ וְלֹא יַשְׁחִיתֶךָ וְלֹא יִשְׁכַּח אֶת בְּרִית אֲבֹתֶיךָ אֲשֶׁר נִשְׁבַּע לָהֶם׃}
}
{
	\eDtr2{
		He won’t let you down, and he won’t destroy you, because
		YHWH, your God, is a merciful God, and He won’t forget
		your fathers’ covenant that He swore to them.
	}
}
{
%	\footnote{
%		a merciful God. This is the first time that Moses mentions to the people
%		a portion of the divine formula that God revealed to him at Sinai in the
%		moment of Moses’ greatest revelation, when he saw God (Exod 34:6). Until
%		this point, Moses had only spoken of it to God Himself (when Moses
%		pleaded with God not to destroy the people in the spies episode; Num
%		14:18). Later these come to be the most quoted words of God in the
%		Tanak: in the prophets (Joel 2:13; Jon 4:2; Nah 1:3), Psalms (86:15;
%		103:8; 111:4; 112:4; 145:8), and later history (2 Chr 30:9; Neh
%		9:17,31). During their travels the people have seen much of God’s just,
%		punishing side. But now, as they are about to begin their life in the
%		promised land, Moses informs them that God is, in the first place, not a
%		wrathful God, but a merciful God.
%	}
}

\Verse{32}
{
	\hDtr1{כִּי שְׁאַל נָא לְיָמִים רִאשֹׁנִים אֲשֶׁר הָיוּ לְפָנֶיךָ לְמִן הַיּוֹם אֲשֶׁר בָּרָא אֱלֹהִים אָדָם עַל הָאָרֶץ וּלְמִקְצֵה הַשָּׁמַיִם וְעַד קְצֵה הַשָּׁמָיִם הֲנִהְיָה כַּדָּבָר הַגָּדוֹל הַזֶּה אוֹ הֲנִשְׁמַע כָּמֹהוּ׃}
}
{
	\eDtr1{
		Because ask of the earliest days that were before you,
		from the day that God created a human on the earth, and
		from one end of the skies to the other end of the skies:
		has there been anything like this great thing? or has
		any-thing like it been heard of?
	}
}
{
}

\Verse{33}
{
	\hDtr1{הֲשָׁמַע עָם קוֹל אֱלֹהִים מְדַבֵּר מִתּוֹךְ הָאֵשׁ כַּאֲשֶׁר שָׁמַעְתָּ אַתָּה וַיֶּחִי׃}
}
{
	\eDtr1{
		Has a people heard God’s voice speaking from inside a fire
		the way you heard—and lived?
	}
}
{
}

\Verse{34}
{
	\hDtr1{אוֹ הֲנִסָּה אֱלֹהִים לָבוֹא לָקַחַת לוֹ גוֹי מִקֶּרֶב גּוֹי בְּמַסֹּת בְּאֹתֹת וּבְמוֹפְתִים וּבְמִלְחָמָה וּבְיָד חֲזָקָה וּבִזְרוֹעַ נְטוּיָה וּבְמוֹרָאִים גְּדֹלִים כְּכֹל אֲשֶׁר עָשָׂה לָכֶם יְהֹוָה אֱלֹהֵיכֶם בְּמִצְרַיִם לְעֵינֶיךָ׃}
}
{
	\eDtr1{
		Or has God put it to the test, to come to take for Himself
		a people from among another people with tests, with signs,
		and with wonders and with war and with a strong hand and
		with an outstretched arm and with great fears like
		everything that YHWH, your God, has done for you in Egypt
		before your eyes?
	}
}
{
%	\footnote{
%		a strong hand and an outstretched arm. More correctly, the first noun
%		(y[image]d) means “forearm” and the second (z[image]r[image]‘a) means
%		“arm.” In Hebrew, k[image]p is the hand, up to the wrist; yãd is the
%		forearm, from the fingers up to the elbow; and z[image]r[image]‘a is the
%		entire arm, from the fingers up to the shoulder (see Gen 49:24). I have
%		retained the traditional translation, “a strong hand and an outstretched
%		arm,” because it is a famous biblical phrase. And I frequently translate
%		y[image]d as hand because the word sometimes can mean just the hand and
%		because of factors of English usage. (See also the comment on Deut
%		11:24.)
%	}
}

\Verse{35}
{
	\hDtr1{אַתָּה הרְאֵתָ לָדַעַת כִּי יְהֹוָה הוּא הָאֱלֹהִים אֵין עוֹד מִלְּבַדּוֹ׃}
}
{
	\eDtr1{
		You have been shown in order to know that YHWH: He is God.
		There is no other outside of Him. 4:35,39. He is God. The
		Hebrew is written with the definite article (literally,
		“He is the God,” not just “He is God” or “He is a God”).
		Grammatically, it excludes the possibility of saying the
		same thing about any other god as well. That is, it is a
		purely monotheistic statement. The sentence that follows
		it in v. 35 confirms this: “There is no other outside of
		Him.” The sentence that follows it in v. 39 does likewise:
		“There isn’t another.” And these sentences were written by
		the seventh century B.C.E. at the very latest. They weigh
		against the claim of some scholars that monotheism was a
		late concept in biblical Israel (later than the Babylonian
		exile in 587 B.C.E.).
	}
}
{
}

\Verse{36}
{
	\hDtr1{מִן הַשָּׁמַיִם הִשְׁמִיעֲךָ אֶת קֹלוֹ לְיַסְּרֶךָּ וְעַל הָאָרֶץ הֶרְאֲךָ אֶת אִשּׁוֹ הַגְּדוֹלָה וּדְבָרָיו שָׁמַעְתָּ מִתּוֹךְ הָאֵשׁ׃}
}
{
	\eDtr1{
		From the skies He had you hear His voice in order to
		discipline you, and on the earth He showed you His great
		fire, and you heard His words from inside the fire.
	}
}
{
}

\Verse{37}
{
	\hDtr1{וְתַחַת כִּי אָהַב אֶת אֲבֹתֶיךָ וַיִּבְחַר בְּזַרְעוֹ אַחֲרָיו וַיּוֹצִאֲךָ בְּפָנָיו בְּכֹחוֹ הַגָּדֹל מִמִּצְרָיִם׃}
}
{
	\eDtr1{
		And because He loved your fathers He chose their seed
		after them, so He brought you out in front of Him from
		Egypt by His great power,
	}
}
{
%	\footnote{
%		their seed after them. Grammatically, the text reads “his seed after
%		him”—using the singular even though the antecedent is the plural “your
%		fathers.” I do not know if there is some special meaning to this.
%	}
}

\Verse{38}
{
	\hDtr1{לְהוֹרִישׁ גּוֹיִם גְּדֹלִים וַעֲצֻמִים מִמְּךָ מִפָּנֶיךָ לַהֲבִיאֲךָ לָתֶת לְךָ אֶת אַרְצָם נַחֲלָה כַּיּוֹם הַזֶּה׃}
}
{
	\eDtr1{
		to dispossess bigger and more powerful nations than you in
		front of you, to bring you, to give you their land as a
		legacy as it is today.
	}
}
{
}

\Verse{39}
{
	\hDtr1{וְיָדַעְתָּ הַיּוֹם וַהֲשֵׁבֹתָ אֶל לְבָבֶךָ כִּי יְהֹוָה הוּא הָאֱלֹהִים בַּשָּׁמַיִם מִמַּעַל וְעַל הָאָרֶץ מִתָּחַת אֵין עוֹד׃}
}
{
	\eDtr1{
		And you shall know today and store it in your heart that
		YHWH: He is God in the skies above and on the earth below.
		There isn’t another.
	}
}
{
}

\Verse{40}
{
	\hDtr1{וְשָׁמַרְתָּ אֶת חֻקָּיו וְאֶת מִצְוֺתָיו אֲשֶׁר אָנֹכִי מְצַוְּךָ הַיּוֹם אֲשֶׁר יִיטַב לְךָ וּלְבָנֶיךָ אַחֲרֶיךָ וּלְמַעַן תַּאֲרִיךְ יָמִים עַל הָאֲדָמָה אֲשֶׁר יְהֹוָה אֱלֹהֶיךָ נֹתֵן לְךָ כּל הַיָּמִים׃}
}
{
	\eDtr1{
		And you shall observe His laws and His commandments that I
		command you today so it will be good for you and for your
		children after you, and so that you’ll extend days on the
		land that YHWH, your God, is giving you forever.”
	}
}
{
}

\Verse{41}
{
	\hDtr1{אָז יַבְדִּיל מֹשֶׁה שָׁלֹשׁ עָרִים בְּעֵבֶר הַיַּרְדֵּן מִזְרְחָה שָׁמֶשׁ׃}
}
{
	\eDtr1{
		Then Moses distinguished three cities across the Jordan
		toward the sun’s rising,
	}
}
{
}

\Verse{42}
{
	\hDtr1{לָנֻס שָׁמָּה רוֹצֵחַ אֲשֶׁר יִרְצַח אֶת רֵעֵהוּ בִּבְלִי דַעַת וְהוּא לֹא שֹׂנֵא לוֹ מִתְּמֹל שִׁלְשֹׁם וְנָס אֶל אַחַת מִן הֶעָרִים הָאֵל וָחָי׃}
}
{
	\eDtr1{
		for the manslayer who would slay his neighbor without
		knowing—and he had not hated him from the day before
		yesterday—to flee there, so he would flee to one of these
		the cities and would live:
	}
}
{
%	\footnote{
%		from the day before yesterday. Meaning: he had no prior hatred of the
%		person he killed. The killing was accidental and without malice.
%	}
}

\Verse{43}
{
	\hDtr1{אֶת בֶּצֶר בַּמִּדְבָּר בְּאֶרֶץ הַמִּישֹׁר לָראוּבֵנִי וְאֶת רָאמֹת בַּגִּלְעָד לַגָּדִי וְאֶת גּוֹלָן בַּבָּשָׁן לַמְנַשִּׁי׃}
}
{
	\eDtr1{
		Bezer in the wilderness of the plain for the Reubenites,
		and Ramoth in Gilead for the Gadites, and Golan in Bashan
		for the Manassites.
	}
}
{
}

\Verse{44}
{
	\hDtr1{וְזֹאת הַתּוֹרָה אֲשֶׁר שָׂם מֹשֶׁה לִפְנֵי בְּנֵי יִשְׂרָאֵל׃}
}
{
	\eDtr1{
		And this is the instruction that Moses set before the
		children of Israel.
	}
}
{
%	\footnote{
%		this is the instruction. In the original context here in Deuteronomy,
%		the word “Torah/instruction” refers only to the things that Moses says
%		to the people across the Jordan from Canaan in his last speech
%		(4:44–46). Jews now sing this verse when the Torah is held up for all to
%		see after it is read each week. These words, “this is the Torah that
%		Moses set … ,” have therefore come to be understood in that context to
%		refer to the entire five books on the Torah scroll.
%	}
}

\Verse{45}
{
	\hDtr1{אֵלֶּה הָעֵדֹת וְהַחֻקִּים וְהַמִּשְׁפָּטִים אֲשֶׁר דִּבֶּר מֹשֶׁה אֶל בְּנֵי יִשְׂרָאֵל בְּצֵאתָם מִמִּצְרָיִם׃}
}
{
	\eDtr1{
		These are the testimonies and the laws and the judgments
		that Moses spoke to the children of Israel when they went
		out from Egypt,
	}
}
{
%	\footnote{
%		testimonies. Meaning: things that have been witnessed. Moses has called
%		the skies and the earth to witness what he has instructed the people
%		(4:26). He will do so again at the end of Deuteronomy (31:28; 32:1).
%	}
}

\Verse{46}
{
	\hDtr1{בְּעֵבֶר הַיַּרְדֵּן בַּגַּיְא מוּל בֵּית פְּעוֹר בְּאֶרֶץ סִיחֹן מֶלֶךְ הָאֱמֹרִי אֲשֶׁר יוֹשֵׁב בְּחֶשְׁבּוֹן אֲשֶׁר הִכָּה מֹשֶׁה וּבְנֵי יִשְׂרָאֵל בְּצֵאתָם מִמִּצְרָיִם׃}
}
{
	\eDtr1{
		across the Jordan in the valley opposite Beth Peor, in the
		land of Sihon, king of the Amorites, who lived in Heshbon,
		whom Moses and the children of Israel struck when they
		came out from the land of Egypt.
	}
}
{
}

\Verse{47}
{
	\hDtr1{וַיִּירְשׁוּ אֶת אַרְצוֹ וְאֶת אֶרֶץ עוֹג מֶלֶךְ הַבָּשָׁן שְׁנֵי מַלְכֵי הָאֱמֹרִי אֲשֶׁר בְּעֵבֶר הַיַּרְדֵּן מִזְרַח שָׁמֶשׁ׃}
}
{
	\eDtr1{
		And they took possession of his land, and of the land of
		Og, king of Bashan, the two kings of the Amorites who were
		across the Jordan at the sun’s rising,
	}
}
{
}

\Verse{48}
{
	\hDtr1{מֵעֲרֹעֵר אֲשֶׁר עַל שְׂפַת נַחַל אַרְנֹן וְעַד הַר שִׂיאֹן הוּא חֶרְמוֹן׃}
}
{
	\eDtr1{
		from Aroer, which is on the edge of the Wadi Arnon, to
		Mount Sion—that is Hermon—
	}
}
{
}

\Verse{49}
{
	\hDtr1{וְכל הָעֲרָבָה עֵבֶר הַיַּרְדֵּן מִזְרָחָה וְעַד יָם הָעֲרָבָה תַּחַת אַשְׁדֹּת הַפִּסְגָּה׃}
}
{
	\eDtr1{
		and all the plain across the Jordan eastward to the sea of
		the plain below the slopes of Pisgah.
	}
}
{
}
